\chapter{The five convolutional autoencoders}\label{ch:convae}

\begin{onething}
Before the transformers, the project built five convolutional autoencoders for one thermal band, Pratibimba,
Antardhana, Tirohita, Asanskrita and Drashta: two classes, one stack of stride-2 convolutions and its mirror, between
264,961 and 1,379,393 trainable numbers, which differ in how they hide pixels, which loss they learn from and whether
they normalise; each step answered a fault of the one before, and you can count every one of them by hand.
\end{onething}

\Cref{ch:recon} listed seven codenames; \cref{ch:shapes} counted one convolutional encoder. This chapter is the
condensed tour of the five convolutional models, one section each: the class, the shapes of one forward pass, the
parameter count, why the model exists, and where it sits in the product. It is shorter than its neighbours on
purpose. These models are not the product's default for anything (Chakshu is, for thermal scenes,
\cref{ch:chakshu}), but they are in the model volume, the scoring dialog offers them, and their history explains the
design of everything after them.

Every count and shape tagged \tagC{} comes from \file{research/textbook/tools/compute_convae.py}, which loads each
model's served checkpoint (copied from the stack's \file{/models/} volume on 2026-10-31) into the real class with a
strict load and records shapes with forward hooks; its output is \file{compute_convae_output.txt}. Paths abbreviate:
\file{components/} is \file{app/foundation_models/components/}, \file{trainers/} and \file{inferencers/} likewise.

\begin{analogy}
Five apprentices learn to restore damaged photographs. The first is only ever shown undamaged photographs and asked
to copy them. The second is shown photographs with a few holes punched in them, the third with most of the picture
punched out. The fourth gets the punched photograph and, on a separate sheet, a stencil showing where the holes are,
so a hole is never mistaken for a black object. The fifth gets the stencil too, and every photograph is first
adjusted to the same brightness. On the job all five are handed photographs with every other dot punched out. Only
the last two were trained for anything like that.
\end{analogy}

%======================================================================
\section{Two building blocks: the stride-2 convolution and its transpose}\label{sec:convae:blocks}

\begin{sloppypar}
All five use \code{SpatialEncoder} and \code{SpatialDecoder} (\file{components/spatial_encoder.py},
\file{components/spatial_decoder.py}). An encoder stage is a convolution with kernel $k=4$, stride $s=2$, padding
$p=1$, then BatchNorm, GELU and a channel dropout of 0.3; by the output-size formula of \cref{ch:shapes}, a side of
$H$ becomes $\lfloor(H+2-4)/2\rfloor+1=H/2$ for even $H$. The channel count starts at a \emph{base} width $c$ and
doubles each stage: $c, 2c, 4c, \dots$
\end{sloppypar}

\begin{definition}[BatchNorm, briefly]
A BatchNorm\index{BatchNorm} layer over $c$ channels standardises each channel with the batch's mean and variance in
training (and with stored running averages at inference), then applies a learned scale and shift per channel: $2c$
parameters, plus $2c+1$ buffers (\cref{ch:compressor} works through it).
\end{definition}

\begin{definition}[Transposed convolution]
A \emph{transposed convolution}\index{transposed convolution} with stride $s$ spreads each input value over the
output: input position $i$ adds $x_i$ times the kernel into the output starting at position $si$, overlapping
contributions add, and $p$ values are cropped from each end. Its output size is
\begin{equation}\label{eq:convae:tconv}
H_\text{out} = (H_\text{in}-1)\,s-2p+k ,
\end{equation}
so with $k=4$, $s=2$, $p=1$: $(H-1)\cdot2-2+4=2H$, an exact doubling. Its parameters count like a convolution's,
$c_\text{in}c_\text{out}k^2+c_\text{out}$.
\end{definition}

\begin{worked}[label={we:convae:tconv}]{A transposed convolution in one dimension}
Input $(1,2)$, kernel $(1,1,1,1)$, stride 2, padding 1, no bias.
\begin{align*}
x_0=1 \text{ at } 0 &: \text{adds } (1,1,1,1) \text{ to positions } 0\text{--}3\\
x_1=2 \text{ at } 2 &: \text{adds } (2,2,2,2) \text{ to positions } 2\text{--}5\\
\text{sum} &: (1,\ 1,\ 1+2,\ 1+2,\ 2,\ 2) = (1,1,3,3,2,2)\\
\text{crop } p=1 \text{ each end} &: (1,3,3,2)
\end{align*}
Two values became four, as \cref{eq:convae:tconv} says: $(2-1)\times2-2+4=4$.
\end{worked}

A decoder stage is the transpose, $k=4$, $s=2$, $p=1$, halving channels, with BatchNorm, GELU and dropout, except
the last, which is the bare transposed convolution. The code's comment explains why: the last layer must emit real
temperatures, which BatchNorm would force to zero mean and GELU would clip below zero
(\file{components/spatial_decoder.py}).

\paragraph{The counting rule for one stage.} With the rules of \cref{ch:shapes},
\begin{equation}\label{eq:convae:stage}
P_\text{enc stage}(a\to b) = 16ab+b+2b = 16ab+3b, \qquad P_\text{dec stage}(a\to b) = 16ab+3b
\ \ (\text{last: } 16ab+b).
\end{equation}

%======================================================================
\section{The five side by side}\label{sec:convae:table}

\begin{table}[htbp]
\centering\small
\begin{tabular}{@{}llllllr@{}}
\toprule
codename & class & encoder channels & hides & loss & norm. & trainable\\
\midrule
Pratibimba & A & $1\to32\to64\to128\to256$ & nothing & MSE & yes & 1,379,393\\
Antardhana & A & $1\to64\to128$ & 13--25\% & MSE & yes & 264,961\\
Tirohita & A & $1\to32\to64\to128$ & 50--75\% & L1 & yes & 329,665\\
Asanskrita & B & $3\to64\to128$ & 35--55\% & L1 & no & 267,009\\
Drashta & C & $3\to64\to128$ & 50--75\% & L1 & yes & 267,009\\
\bottomrule
\end{tabular}
\caption{The five convolutional models as served. Classes: A \texttt{SpatialAutoencoder}, B
\texttt{UnNormalizedSpatialAutoencoder}, C \texttt{NormalizedMaskedSpatialAutoencoder}. ``Hides'' is the share of valid
pixels hidden per training patch, drawn uniformly in the range for each patch. Counts from strict loads of the served checkpoints \tagC; masking ranges
and losses from \texttt{trainers/}. Asanskrita's and Drashta's encoders read three channels: the picture, a validity
mask and an input mask.}
\label{tab:convae:five}
\end{table}

All five were trained on Landsat~9 thermal patches in °C at sizes 128, 256 and 512 within each epoch, with Adam at
$10^{-3}$, and all five discard patches with under 40\% valid pixels (\file{trainers/}, \code{MIN_VALID_PIXEL_FRACTION}).
The four normalising models carry one mean and one spread for Landsat temperatures, $24.5756\pm13.5744$ °C
(\file{app/constants/thermal_pixel_consts.json}); \cref{ch:chakshu}, which uses the same two numbers, works with them
in detail. Every one is served with patch 128, stride 64, and two-pass \emph{pixel} checkerboard
inference: one pass hides every pixel whose row plus column is odd, the other those where it is even
(\file{inferencers/spatial_autoencoder_inferencer.py} and its two siblings).

%======================================================================
\section{Pratibimba: the plain autoencoder}\label{sec:convae:pratibimba}

\paragraph{What it is.} Class A with one input channel, base 32 and 4 stages
(\file{components/spatial_auto_encoder.py}), checkpoint v8.1.0, epoch 29. Codename: ``reflection''.

\begin{center}\small
\begin{tabular}{@{}lll@{}}
\toprule
step & layer & shape ($B=1$, $128\times128$)\\
\midrule
input, normalised & & \shape{(1, 1, 128, 128)}\\
encoder 0--3 & conv $k4\,s2\,p1$, BN, GELU, dropout & \shape{(1, 32, 64, 64)}, \shape{(1, 64, 32, 32)},\\
 & & \shape{(1, 128, 16, 16)}, \shape{(1, 256, 8, 8)}\\
decoder 0--2 & transposed conv, BN, GELU, dropout & \shape{(1, 128, 16, 16)}, \shape{(1, 64, 32, 32)},\\
 & & \shape{(1, 32, 64, 64)}\\
decoder 3 & transposed conv only & \shape{(1, 1, 128, 128)}, then denormalised\\
\bottomrule
\end{tabular}
\end{center}

\paragraph{Count.} Encoder, by \cref{eq:convae:stage}:
$(16\cdot32+96)+(16\cdot32\cdot64+192)+(16\cdot64\cdot128+384)+(16\cdot128\cdot256+768)
=608+32{,}960+131{,}456+525{,}056=690{,}080$. Decoder: $(16\cdot256\cdot128+384)+(16\cdot128\cdot64+192)+(16\cdot64\cdot32+96)
+(16\cdot32+1)=524{,}672+131{,}264+32{,}864+513=689{,}313$. Total 1,379,393 \tagC. The manifest's 1,380,812 adds 1,419
buffers: a running mean and variance for each of the $480+224=704$ BatchNorm channels ($2\times704=1{,}408$), 7 batch
counters, and 4 normalisation values; $1{,}408+7+4=1{,}419$.

\paragraph{Why it exists, and its fault.} It is the baseline: reproduce the patch through a bottleneck. Its trainer
(\file{trainers/spatial_autoencoder_trainer.py}) hides nothing; the loss is the MSE over all valid pixels. That is
the identity trap of \cref{ch:recon}: the only thing stopping a copy is the bottleneck, $256\times8\times8=16{,}384$
values for a 16,384-pixel patch, exactly as many numbers as the input. At inference the product then hides half the
pixels in a checkerboard, a regime the model never saw.

%======================================================================
\section{Antardhana and Tirohita: hide pixels, keep the class}\label{sec:convae:masked}

\paragraph{What they are.} The same \code{SpatialAutoencoder} class. Antardhana (``concealment''), v0.1.0 epoch
121, is base 64 with 2 stages; Tirohita (``the hidden one''), v0.4.0 epoch 77, base 32 with 3 stages. Their trainers
(\file{trainers/spatial_masked_autoencoder_trainer.py}, \file{..._trainer_l1_loss.py}) draw a hiding share per patch,
uniformly in $[0.13, 0.25]$ for Antardhana and $[0.50, 0.75]$ for Tirohita, hide that share of the valid pixels by
setting them to 0, and grade only the hidden ones, with MSE and L1 respectively.

\begin{center}\small
\begin{tabular}{@{}lll@{}}
\toprule
 & Antardhana & Tirohita\\
\midrule
encoder outputs & \shape{(1, 64, 64, 64)} & \shape{(1, 32, 64, 64)}\\
 & \shape{(1, 128, 32, 32)} & \shape{(1, 64, 32, 32)}\\
 & & \shape{(1, 128, 16, 16)}\\
decoder outputs & \shape{(1, 64, 64, 64)} & \shape{(1, 64, 32, 32)}\\
 & \shape{(1, 1, 128, 128)} & \shape{(1, 32, 64, 64)}\\
 & & \shape{(1, 1, 128, 128)}\\
bottleneck values & $128\times32\times32=131{,}072$ & $128\times16\times16=32{,}768$\\
trainable & 264,961 & 329,665\\
\bottomrule
\end{tabular}
\end{center}
Antardhana's stages: $16\cdot64+192=1{,}216$; $16\cdot64\cdot128+384=131{,}456$; decoder $16\cdot128\cdot64+192=131{,}264$
and the last $16\cdot64+1=1{,}025$. Tirohita's encoder is the 165,024 of \cref{we:shapes:spatialcount}; its decoder is
$131{,}264+32{,}864+513=164{,}641$ \tagC.

\paragraph{Why they exist, and their fault.} Hiding pixels removes the identity answer for the hidden ones
(\cref{ch:recon}); Tirohita adds the L1 loss, which does not pull predictions towards anomalies
(\cref{sec:recon:l1l2}), and much heavier hiding. The fault both keep: a hidden pixel is set to 0, then normalised to
$(0-24.5756)/13.5744=-1.81$, which is exactly what a real pixel at 0\,°C looks like. With one input channel the network
cannot tell ``hidden'' from ``freezing''.

\begin{sourcenote}[: two siblings that differ in more than the loss, and their validation losses]
The manifest and \file{components/spatial_auto_encoder.py} present Antardhana and Tirohita as one model with two losses.
They also differ in size (base 64, 2 stages against base 32, 3 stages) and in hiding share (13--25\% against 50--75\%)
(\cref{ch:losses} noted the size difference). The manifest's note that Tirohita has ``Lower val loss than the MSE
sibling'' compares 1.79 (an L1, in °C) with 17.29 (an MSE, in °C$^2$), and both trainers validate with
\emph{nothing hidden} (\lstinline!model(pixels, mask=mask)!), so neither number measures the hidden-pixel task they were
trained on. The two numbers cannot be compared. Also, both manifests name \code{MaskedSpatialAutoencoderInferencer};
the resolver routes both through \code{SpatialAutoencoderInferencer} (their \code{foundation_model_name} is
\code{spatial_autoencoder}), and \file{docs/04-models.md}'s ``two architectures can be trained but not run'' is true of
the enum names only: the product runs both models under Pratibimba's inferencer.
\end{sourcenote}

%======================================================================
\section{Asanskrita: tell the network where the holes are}\label{sec:convae:asanskrita}

\paragraph{What it is.} Class B, base 64, 2 stages
(\file{components/unnormalized_spatial_auto_encoder.py}), v0.3.0 epoch 82. Codename: ``unrefined'', a pun on
\emph{unnormalised}: it is the only model with no normalisation buffers. Its encoder reads three channels:
\begin{equation}\label{eq:convae:three}
\big[\ T\cdot m_\text{valid}\cdot m_\text{input},\ \ m_\text{valid},\ \ m_\text{input}\ \big],
\end{equation}
the temperature in °C with hidden and invalid pixels zeroed, the validity mask, and the input mask (1 where a
pixel is valid and visible). A hidden pixel has $m_\text{valid}=1$, $m_\text{input}=0$; an invalid one has both 0; so
the network can tell ``hidden'' from ``nothing here'' from ``a real 0\,°C''.

\begin{center}\small
\begin{tabular}{@{}lll@{}}
\toprule
step & layer & shape\\
\midrule
concatenate & \eqref{eq:convae:three} & \shape{(1, 3, 128, 128)}\\
encoder 0, 1 & conv $3\to64$, $64\to128$ & \shape{(1, 64, 64, 64)}, \shape{(1, 128, 32, 32)}\\
decoder 0, 1 & transposed $128\to64$, $64\to1$ & \shape{(1, 64, 64, 64)}, \shape{(1, 1, 128, 128)}, in °C\\
\bottomrule
\end{tabular}
\end{center}
Count: the first stage is $16\cdot3\cdot64+64+128=3{,}072+192=3{,}264$; the rest is Antardhana's
($131{,}456+131{,}264+1{,}025$), so $3{,}264+131{,}456+131{,}264+1{,}025=267{,}009$ \tagC. The two extra input channels cost
$3{,}264-1{,}216=2{,}048=2\times16\times64$ weights.

\paragraph{Why, and the trade-off.} The mask channels fix the ``hidden or freezing'' ambiguity. Dropping normalisation
was the experiment: raw °C goes straight into the first convolution, followed by BatchNorm, which re-centres the
channels anyway. Its trainer hides 35--55\% of valid pixels with L1; the configuration's \code{masking_range} is
ignored by this trainer, whose function defaults happen to be the same $[0.35, 0.55]$.

%======================================================================
\section{Drashta: the masks and the normalisation}\label{sec:convae:drashta}

\paragraph{What it is.} Class C, base 64, 2 stages, in
\file{components/normalized_masked_spatial_auto_encoder.py}; v1.2.0 epoch 69. Codename: ``the seer''. Asanskrita's
architecture with \code{PixelNormalize} on the way in and \code{PixelDenormalize} on the way out; its trainer reads
\code{masking_range} $[0.5, 0.75]$ from the configuration, with L1. Shapes and count are Asanskrita's: 267,009
trainable, and 519 buffers against Asanskrita's 515 (the four normalisation values) \tagC.

\begin{sourcenote}[: a hidden pixel is $-1.81$, not 0]
The class's docstring says channel 0 holds ``normalized thermal pixels (zeroed where input\_mask == 0)''. The forward
pass multiplies by the masks \emph{before} normalising, so a hidden pixel enters as $(0-24.5756)/13.5744=-1.81$, the
same value as in Antardhana, and only the mask channels say it is hidden. Harmless, because the mask channels are
there; but read the docstring as describing the raw picture, not channel 0.
\end{sourcenote}

%======================================================================
\section{What the product does with them}\label{sec:convae:product}

\begin{center}\small
\begin{tabular}{@{}llll@{}}
\toprule
codename & inferencer (as routed) & default score & normalisation file\\
\midrule
Pratibimba & \texttt{SpatialAutoencoderInferencer} & MSE & thermal\\
Antardhana & \texttt{SpatialAutoencoderInferencer} & MSE & thermal\\
Tirohita & \texttt{SpatialAutoencoderInferencer} & L1 & thermal\\
Asanskrita & \texttt{MaskedSpatialAutoencoderInferencer} & L1 & none\\
Drashta & \texttt{NormalizedMaskedAutoencoderInferencer} & L1 & thermal\\
\bottomrule
\end{tabular}
\end{center}
All five: scoring methods L1 and MSE, patch 128, stride 64, batch 8, sensor family thermal
(\file{backend/allotrope/foundation_models/resolver.py}). The resolver rebuilds each configuration from the manifest's
\code{encoder_dims} string (\code{"1→32→64→128"}: base 32, three stages; \code{"3→64→128"}: input 3, so constructor
\code{in_channels} $3-2=1$), and every rebuilt configuration matches its checkpoint \tagC.

Three behaviours differ from Chakshu's inferencer (\cref{ch:chakshu}): the masking is per pixel, not per token; the
validity mask is not eroded; and a pixel no tile covers is set to 0, not to its input value, so it scores its whole
temperature as residual (\file{docs/09-known-issues.md}, item 4).

\paragraph{A synthetic comparison.} The script runs each model, through the product's own inferencer and with its
served weights, on a smooth synthetic field, $T=25+3\sin(x/40)+2\cos(y/30)$\,°C for column $x$ and row $y$ of a
$128\times128$ patch, and on copies with one hot pixel or a $3\times3$ hot spot at 45\,°C planted near the centre
\tagA:
\begin{center}\small
\begin{tabular}{@{}lccc@{}}
\toprule
model & background (median) & one hot pixel & $3\times3$ spot, centre\\
\midrule
Pratibimba & 10.87\,°C & 19.20 & 17.60\\
Antardhana & 6.59 & 24.23 & 22.05\\
Tirohita & 2.32 & 20.21 & 14.78\\
Asanskrita & 0.75 & 19.37 & 16.57\\
Drashta & 0.58 & 18.14 & 16.42\\
\midrule
Chakshu ($256\times256$) & 0.019 & 20.96 & 7.59\\
\bottomrule
\end{tabular}
\end{center}
All values are $|T-\hat T|$ in °C \tagC; Chakshu's row is from \cref{ch:chakshu}, which runs the same field at its own
patch size. On a field this smooth the background error is the number to read: a detector's hot pixel must stand above
it. Pratibimba's 10.9\,°C, half the planted anomaly, is what happens when a model trained on whole pictures meets a
checkerboard of holes it has never seen; the three single-channel models improve in step with how much hiding they
trained on; the two mask-channel models are an order of magnitude better again; Chakshu is thirty times better
than the best of them. One synthetic field is not a benchmark: nothing here was measured on a labelled thermal scene,
and none exists in the repository.

\begin{sourcenote}[: the per-scene override and these models]
Four of the five carry \code{normalize.*} buffers in their checkpoints. \Cref{sec:chakshu:hotsat}, later, shows that loading a
checkpoint writes those buffers back over a per-scene override, so the HotSat-1 plan fails for these four exactly as
for Chakshu. Asanskrita has no normalisation to override: it would read HotSat's digital numbers as degrees.
\end{sourcenote}

\begin{drift}
Read 2026-10-31 from the \file{/models/} volume's manifests: Pratibimba v8.1.0 (epoch 29, validation 13.77, MSE),
Antardhana v0.1.0 (epoch 121, 17.29, MSE), Tirohita v0.4.0 (epoch 77, 1.79, L1), Asanskrita v0.3.0 (epoch 82, 2.15, L1),
Drashta v1.2.0 (epoch 69, 5.53, L1). Manifest parameter figures (1,380,812; 265,480; 330,314; 267,524; 267,528) include
buffers. Validation for every masked model is computed with nothing hidden.
\end{drift}

%======================================================================
\begin{confusions}
\confusion{Antardhana and Tirohita differ only in the loss.}{They also differ in width, depth and how much they hide
(13--25\% against 50--75\%).}
\confusion{Tirohita's validation loss is ten times lower, so it is ten times better.}{One is an L1 in °C, the other an
MSE in °C$^2$, and both are measured with nothing hidden. They are not comparable.}
\confusion{The masked models cannot be run: no inferencer is registered for them.}{The resolver routes Antardhana and
Tirohita through Pratibimba's inferencer; only the enum names lack an entry.}
\confusion{Asanskrita is unnormalised, so its network sees raw temperatures throughout.}{Only its first convolution
does; the BatchNorm after it standardises every channel, with learned scale and shift.}
\confusion{A hidden pixel in Drashta is a zero.}{It is $-1.81$ after normalisation; the mask channels mark it.}
\confusion{More parameters, better model.}{Pratibimba has five times Drashta's parameters and an error nineteen times
larger on the synthetic field ($10.87/0.58=18.7$). What it was trained to do matters more.}
\end{confusions}

\begin{checkyourself}
\item\label{ex:convae:1} Use \cref{eq:convae:tconv} to find the output length of a transposed convolution with $k=4$,
$s=2$, $p=1$ on an input of 16, and with $k=3$, $s=2$, $p=1$ on 16. Which one doubles?
\item\label{ex:convae:2} Count by hand a \code{SpatialAutoencoder} with base 16 and 3 stages (one input channel), encoder
and decoder separately.
\item\label{ex:convae:3} How many values does Pratibimba's bottleneck hold for a $256\times256$ patch, and how does that
compare with the input?
\item\label{ex:convae:4} A pixel is valid and hidden. Write its three input channels for Asanskrita and for Drashta, if
its true temperature is 30\,°C.
\item\label{ex:convae:5} Explain in two sentences why Pratibimba's background error on the synthetic field is so large,
using the words ``trained'' and ``checkerboard''.
\item\label{ex:convae:6} Repeat \cref{we:convae:tconv} with input $(3,-1)$ and kernel $(1,2,2,1)$.
\end{checkyourself}

\nextchapter{ch:compressor}
